\begin{infoclean}
    \centering
    Seja $x$ um número binário de 4 bits, $x_i$ um índice $i$ de $x$ e $A,B,C,D$ resultados temporários, o cálculo de $x^2$ é dado por:
\end{infoclean}
\vspace{-.5cm}
\begin{gather}
    x^2 = \begin{cases}
        \text{com $i=3 \rightarrow$ Se $x_i = 1 \Rightarrow A= LSL_{x}^i$ caso contrário $A=0$} \\
        \text{com $i=2 \rightarrow$ Se $x_i = 1 \Rightarrow B=LSL_{x}^i$ caso contrário $B=0$}  \\
        \text{com $i=1 \rightarrow$ Se $x_i = 1 \Rightarrow C=LSL_{x}^i$ caso contrário $C=0$}  \\
        \text{com $i=0 \rightarrow$ Se $x_i = 1 \Rightarrow D=LSL_{x}^i$ caso contrário $D=0$}
    \end{cases} \\
    x^2 = A + B + C + D
\end{gather}